\documentclass[11pt,reqno]{amsart}

\usepackage{amsmath,amssymb,amsthm}
\usepackage{mathtools}
\usepackage[T1]{fontenc}
\usepackage{lmodern}
\usepackage{microtype}
\usepackage{stmaryrd} % \llbracket \rrbracket
\usepackage[colorlinks=true,linkcolor=blue!55!black,citecolor=blue!55!black,urlcolor=blue!55!black]{hyperref}
\usepackage{tikz-cd}

% ---- theorem environments ----
\theoremstyle{plain}
\newtheorem{theorem}{Theorem}[section]
\newtheorem{proposition}[theorem]{Proposition}
\newtheorem{lemma}[theorem]{Lemma}
\newtheorem{corollary}[theorem]{Corollary}
\newtheorem{conjecture}[theorem]{Conjecture}

\theoremstyle{definition}
\newtheorem{definition}[theorem]{Definition}
\newtheorem{example}[theorem]{Example}
\newtheorem{remark}[theorem]{Remark}
\newtheorem{schema}[theorem]{Schema}

% ---- macros ----
\newcommand{\Set}{\mathrm{Set}}
\newcommand{\Poly}{\mathrm{Poly}}
\newcommand{\Cat}{\mathrm{Cat}}
\newcommand{\Cont}{\mathrm{Cont}}
\newcommand{\DCont}{\mathrm{DCont}}
\newcommand{\Fam}{\mathrm{Fam}}
\newcommand{\Struct}{\mathrm{Struct}}
\newcommand{\EM}{\mathrm{EM}}
\newcommand{\ob}{\mathrm{ob}}
\newcommand{\cod}{\mathrm{cod}}
\newcommand{\dom}{\mathrm{dom}}
\newcommand{\id}{\mathrm{id}}
\newcommand{\leaves}{\mathrm{leaves}}
\newcommand{\graft}{\mathrm{graft}}
\newcommand{\spl}{\mathrm{split}}
\newcommand{\lf}{\mathsf{lf}}
\newcommand{\nd}{\mathsf{nd}}
\newcommand{\EF}{\mathbb{E}}
\newcommand{\Peb}{\mathbb{P}}
\newcommand{\Mod}{\mathbb{M}}
\newcommand{\sem}[1]{\llbracket #1 \rrbracket}
\newcommand{\pos}{P}
\newcommand{\Sh}{S}
\newcommand{\cext}[1]{\sem{#1}}
\newcommand{\pathfun}{\mathsf{p}}
\newcommand{\A}{\mathcal{A}}
\newcommand{\T}{\mathcal{T}}
\newcommand{\Paths}{\mathcal{P}}
\newcommand{\C}{\mathcal{C}}
\newcommand{\op}{\mathrm{op}}
\newcommand{\defeq}{\mathrel{\vcentcolon=}}
\DeclareMathOperator{\Sub}{Sub}

\begin{document}

\title[The rooted-tree normal form]{The rooted-tree normal form:\\
directed containers, free monoids, and arboreal game comonads}

\author{MacBeth}
\address{Kodamai}
\email{neil@kodamai.com}

\date{2 October 2026}

\begin{abstract}
A recurring move in the theory of containers and polynomial functors replaces an
\emph{arbitrary-position} condition by a \emph{rooted, tree-shaped} discipline, under which the
operative backward map reads only a unique root-path. We isolate this move and show that it is the
same move in four independently-developed settings: the root axiom that turns a container into a
comonad (Ahman--Uustalu), the path-confinement of the backward pass of a branching free monoid, the
tree-connectedness axiom that repairs the theory of arboreal categories (Jakl--Reggio), and the
forest condition on polynomial comonoids. The move has a single geometric content---in a forest the
path to the root is unique---which in the last and sharpest instance becomes a theorem: the
Eilenberg--Moore coalgebras of a polynomial comonad $c = C$ are the copresheaves on $C$, so the
coalgebra category is arboreal exactly when $C$ is a forest. This identifies the arboreal coalgebra
categories of finite model theory with copresheaf categories of polynomial comonads, giving a bridge
between the game-comonad programme and the Poly programme. We close with an honestly-graded open
problem on how the two meet at the level of fibrations.
\end{abstract}

\maketitle

\section{Introduction}

\subsection*{A move one keeps making}
Here is a pattern. One is handed a construction---a container, a free monoid, a category of
coalgebras---together with a condition phrased over \emph{all} positions, or \emph{all}
coproducts, or \emph{all} arrows. The condition looks hard to control because the positions branch:
different answers lead to different continuations, different subobjects combine in uncontrolled ways.
Then one notices that the condition can be re-phrased so that the data is confined to a
\emph{rooted tree}, and that the map one actually has to understand reads only a \emph{single
root-path}---a unique prefix, a unique minimal element. The branching that threatened to create an
obstruction is forbidden by the tree shape, and the obstruction evaporates.

This paper is about that move. We have met it four times, in four settings developed independently of
one another, by different communities, for different reasons. The purpose of the paper is to state the
move once, precisely, and to exhibit the four settings as instances of it---so that the reader who
knows one of them recognises the others, and so that a tool sharpened in one can be carried to the
rest.

\subsection*{The four settings}
The four instances, in order of increasing categorical sharpness, are these.

\begin{enumerate}
\item[(I)] \emph{The root axiom.} A container $(S \triangleleft P)$ is a comonad on $\Set$---equivalently
a small category---exactly when its positions carry a distinguished root at each shape together with a
compatible codomain-and-composition discipline. This is the founding instance: the move is already the
definition of a \emph{directed} container \cite{AhmanChapmanUustalu2014,AhmanUustalu2016}.

\item[(II)] \emph{The prefix cut.} The free monoid on a container $C$ under composition---a type of
answer-driven interaction tree---carries a multiplication whose forward part grafts trees and whose
backward part splits a leaf-path. We prove (Lemma~\ref{lem:D1}) that the backward part is
\emph{independent of the grafted data}: it cuts a leaf-path at its unique prefix that is a leaf of the
outer tree. The feared obstruction to a bidirectional universal property---that answer-dependence
corrupts the backward pass---cannot occur, because the cut sees only the outer tree.

\item[(III)] \emph{Factoring through the least element.} In the theory of \emph{arboreal categories}
\cite{AbramskyShah2021,JaklReggio2026}, an axiom requiring paths to be \emph{connected} fails for the
modal game comonads. Jakl and Reggio repair it with \emph{tree-connectedness}: an arrow into a colimit
of a tree-diagram of paths must factor through a \emph{least} element of the tree
\cite[Def.~11]{JaklReggio2026}. This is the same move---the factoring is through a unique minimal
element---and it recovers all the pebbling and Ehrenfeucht--Fra\"iss\'e results while admitting the
modal comonads.

\item[(IV)] \emph{The base is a forest.} The move becomes a property of a small category. We show
(Lemma~\ref{lem:coalg}) that the Eilenberg--Moore coalgebras of a polynomial comonad $c = C$ are the
copresheaves $[C,\Set]$; for the Ehrenfeucht--Fra\"iss\'e comonad this is the category of forests of
bounded height. Hence (Corollary~\ref{cor:arboreal-forest}) ``\emph{the coalgebra category is arboreal}''
is exactly ``\emph{$C$ is a forest}''. The rooted-tree discipline is now a property of the base
category, and arboreality is that property read through the equivalence \emph{polynomial comonad $=$
directed container $=$ small category}.
\end{enumerate}

\subsection*{The main result, and what is and is not new}
The load-bearing theorem is instance (IV):

\begin{theorem}[see Lemma~\ref{lem:coalg} and Corollary~\ref{cor:arboreal-forest}]\label{thm:main}
Let $c = C$ be a polynomial comonad, viewed as a small category. The category of Eilenberg--Moore
coalgebras of the comonad $\cext{c}$ on $\Set$ is equivalent to the copresheaf category $[C,\Set]$,
i.e.\ to the discrete opfibrations over $C$. When $C$ is a forest category these copresheaves are
forests; for the Ehrenfeucht--Fra\"iss\'e comonad, whose comonoid is the linear order $[k]$, they are
exactly the forests of height $\le k$ of \textup{finite model theory}, and for the pebbling comonad the
trie-indexed forests. The arboreal coalgebra categories of these game comonads are thus the copresheaf
categories of their polynomial comonoids.
\end{theorem}

We are careful about novelty. That the coalgebras of a polynomial comonad are copresheaves on its
category is \emph{folklore-adjacent} in the Poly community---it is the familial / cofree-coalgebra
story \cite{SpivakGarnerFairbanks2025}---and that the Ehrenfeucht--Fra\"iss\'e coalgebras are forests is
proved directly by Abramsky and Shah \cite{AbramskyShah2021}. What we contribute is not a new deep
theorem but the \emph{clean identification tying the two}: the one-line reason that turns ``arboreality
$=$ the base is a forest'' from an analogy into a statement with a proof, and the observation that this
is the fourth and sharpest appearance of a single structural move. The value here is the
\emph{bridge}---between the game-comonad programme of finite model theory and the Poly programme of
compositional systems---not a priority claim.

\subsection*{Why a reader should care}
Two reasons. First, pedagogical: once one sees that the root axiom, the prefix cut, the
tree-connectedness repair, and the forest condition are one move, each becomes easier to remember and
to apply, and a technique from one setting (for instance, the path-confinement argument of
Section~\ref{sec:prefix}) suggests itself in the others. Second, structural: Theorem~\ref{thm:main}
says the arboreal coalgebra categories that finite model theory studies are exactly copresheaf
categories of polynomial comonads. The two programmes have never cited one another on this point; the
identification lets results and intuitions cross. We make the bridge precise, and in
Section~\ref{sec:open} we push it to the level of fibrations, where it becomes a genuine, and genuinely
open, question---stated there with its grade and its obstructions named.

\subsection*{Outline}
Section~\ref{sec:prelim} fixes containers, the composition product, the equivalence of polynomial
comonoids with small categories, directed containers, and forests. Section~\ref{sec:schema} states the
normal form once, as a schema, after isolating its single geometric content. Section~\ref{sec:instances}
walks the four instances in order of increasing sharpness, proving instances (II) and (IV) and citing
instances (I) and (III). Section~\ref{sec:open} reads the four as one and states the open problem that
the schema opens, honestly graded. Section~\ref{sec:conclusion} concludes.

\medskip
\noindent\emph{Provenance.} Instances (I) and (III) are established and cited; instances (II) and (IV)
are proved here, the former by structural induction with a computational check, the latter by a direct
argument. The material of Section~\ref{sec:open} beyond the stated theorem is explicitly graded
\emph{computed} or \emph{open}. We have tried to keep the claims exactly as strong as the proofs behind
them, and no stronger.

\section{Preliminaries}\label{sec:prelim}

We recall only what the four instances use. Throughout, $\Set$ is the category of sets.

\subsection{Containers and polynomial functors}
A \emph{container} $C = (S \triangleleft P)$ is a set $S$ of \emph{shapes} together with a family
$P : S \to \Set$ of \emph{positions}. Its \emph{extension} is the functor
\[
  \cext{C} : \Set \to \Set, \qquad \cext{C} A = \sum_{s \in S} A^{P s},
\]
a coproduct of representables; functors of this form are the \emph{polynomial} functors. A morphism of
containers $\varphi : (A \triangleleft P_A) \to (B \triangleleft P_B)$ is a \emph{forward} shape map
$\varphi_1 : A \to B$ together with a \emph{backward} family
$\varphi^\sharp_a : P_B(\varphi_1 a) \to P_A(a)$, contravariant in positions. Containers and their
morphisms form a category $\Cont \cong \Fam(\Set^{\op})$, often written $\Poly$. We use the two freely.

The reading that organises instance (II) is the \emph{bidirectional} one: a shape $s \in S$ is a
request, and a position $p \in P s$ is a possible answer to it.

\begin{proposition}[Gambino--Kock \cite{GambinoKock2013}, Thm.~4.5]\label{prop:ff}
The extension functor $\cext{-} : \Poly \to [\Set,\Set]$ is fully faithful: for containers $p,q$,
\[
  \Poly(p,q) \;\cong\; \mathrm{Nat}(\cext{p},\cext{q}).
\]
\end{proposition}

\subsection{The composition product and polynomial comonoids}
The \emph{composition product} $\triangleleft$ on $\Poly$ computes functor composition,
$\cext{p \triangleleft q} = \cext{p} \circ \cext{q}$, with unit $y = (1 \triangleleft 1)$, the identity
functor. Concretely,
\[
  (G \triangleleft F).\mathrm{Shape} = \sum_{t \in T}\big(Q t \to S\big), \qquad
  (G \triangleleft F).\mathrm{Pos}(t,f) = \sum_{q \in Q t} P(f q).
\]
A \emph{comonoid} in the monoidal category $(\Poly, \triangleleft, y)$ is a container $c$ with a counit
$\varepsilon : c \to y$ and a comultiplication $\delta : c \to c \triangleleft c$ satisfying the
comonoid laws.

\begin{proposition}[Ahman--Uustalu \cite{AhmanUustalu2016}; Spivak--Garner--Fairbanks
\cite{SpivakGarnerFairbanks2025}]\label{prop:comonoid-cat}
Comonoids in $(\Poly,\triangleleft,y)$ are precisely small categories. Under the equivalence, a comonoid
$c = (S \triangleleft P)$ has objects $S$ and arrows-out-of-$s$ the positions $P s$; the counit picks, at
each $s$, an identity arrow $\id_s \in P s$; the comultiplication provides, for each $s$, a codomain map
$\cod_s : P s \to S$ and a composition $\sum_{f \in P s} P(\cod_s f) \to P s$, $(f,g) \mapsto g \circ f$.
The comonoid laws are the category axioms.
\end{proposition}

Combining Propositions~\ref{prop:ff} and~\ref{prop:comonoid-cat}: a comonad on $\Set$ whose underlying
functor is polynomial is \emph{automatically} a polynomial comonoid (the counit and comultiplication,
being natural transformations between polynomial functors, are extensions of unique container morphisms
by full faithfulness, and the comonad laws transport), hence a small category. We write $c = C$ for a
polynomial comonad and the small category it is.

\subsection{Directed containers}
A \emph{directed container} \cite{AhmanChapmanUustalu2014,AhmanUustalu2016} is exactly such a structure
presented element-wise: a container $(S \triangleleft P)$ with, for each shape $s$, a \emph{root}
position $o_s \in P s$, a \emph{subshape} $s \downarrow p \in S$ for each $p \in P s$, and an addition
$p \oplus q \in P s$ for $q \in P(s \downarrow p)$, satisfying unit and associativity laws. These are the
category axioms of Proposition~\ref{prop:comonoid-cat}, with $o_s = \id_s$, $s \downarrow p = \cod_s p$,
and $\oplus$ composition. We write $\DCont \cong \Cat$ for the resulting equivalence.

\subsection{Forests, paths, and copresheaves}
A poset is a \emph{forest} if the down-set $\{x : x \le y\}$ of every element $y$ is a finite chain;
a \emph{tree} is a forest with a least element. We view a small category $C$ as generating a reachability
preorder; $C$ is a \emph{forest category} when that preorder is a forest, so that from every object there
is a \emph{unique} directed path to a root. The \emph{linear order} $[k] = (\{1,\dots,k\},\le)$ and the
$k$-ary \emph{prefix trie} on nonempty words $[k]^{+}$ are forests.

For a small category $C$, a \emph{copresheaf} is a functor $X : C \to \Set$; the copresheaf category
$[C,\Set]$ is equivalently the category of \emph{discrete opfibrations} over $C$. This is the shape in
which instance (IV) is stated.

\section{The normal form}\label{sec:schema}

\subsection{The single geometric fact}
All four instances draw on one fact about forests, which we isolate before stating the schema.

\begin{lemma}[Root-path uniqueness]\label{lem:rootpath}
Let $F$ be a forest. The following coincide for an element $y \in F$:
\begin{enumerate}
\item the down-set $\{x : x \le y\}$, which is a finite chain;
\item the unique directed path from $y$ to its root;
\item the unique \emph{minimal} (equivalently \emph{least}) element of that down-set;
\item for the tree of \emph{prefixes} of a path ending at $y$, the unique prefix lying in any fixed
prefix-closed subset.
\end{enumerate}
In each guise the data attached to $y$ ``below'' it is a chain, carrying no branching.
\end{lemma}

\begin{proof}
A forest is exactly a poset in which every principal down-set is a chain, which is the equivalence of
(1) with the uniqueness of the descending path (2); a nonempty finite chain has a least element, giving
(3); and in a well-founded tree the leaves form a prefix code---no leaf is a proper prefix of another---so
the prefix of a given path lying in a prefix-closed set is unique, giving (4). The last clause restates
that a chain does not branch.
\end{proof}

The four instances differ only in \emph{which} guise of Lemma~\ref{lem:rootpath} they use: instance (I)
distinguishes the root (3); instance (II) cuts at the unique prefix (4); instance (III) factors through
the least element (3); instance (IV) makes ``every down-set is a chain'' (1) a property of the base
category.

\subsection{The schema}
We can now say what the move is.

\begin{schema}[The rooted-tree normal form]\label{sch:main}
A condition phrased over \emph{arbitrary} positions, coproducts, or arrows is replaced by a
\emph{rooted, tree-shaped} discipline: the data is confined to a forest, and the operative backward or
factoring map is determined by the \emph{unique root-path} of Lemma~\ref{lem:rootpath}---a unique
prefix, a least element, a chain down-set. Because a chain does not branch, the branching data that
would have obstructed the condition is never read, and the condition holds.
\end{schema}

Schema~\ref{sch:main} is not itself a theorem---it is the shape that the four theorems of the next
section share. In instance (IV) it becomes a theorem verbatim: the condition (arboreality) is
\emph{equivalent} to the forest property of the base. We take that equivalence as the evidence that the
schema is real and not merely a family resemblance.

\section{The four instances}\label{sec:instances}

\subsection{Instance (I): positions carry a root}\label{sec:root}
The founding instance is the definition of a directed container. A bare container $(S \triangleleft P)$
has \emph{arbitrary} positions; there is no way to compose an observation at $s$ with an observation at a
position of $s$, because a position carries no shape. The directed-container structure supplies exactly
the missing rooted discipline: a root $o_s \in P s$ to serve as identity, a subshape $s \downarrow p$ to
serve as codomain, and an addition to serve as composition (Proposition~\ref{prop:comonoid-cat}). This is
the move at its most elementary---guise (3) of Lemma~\ref{lem:rootpath}, the distinguished least
element---and it is \emph{necessary}: a container is a comonad if and only if it admits such a structure
\cite{AhmanChapmanUustalu2014}, i.e.\ if and only if it is a small category \cite{AhmanUustalu2016}. We
record it as the base case and move on.

\subsection{Instance (II): the prefix cut}\label{sec:prefix}
The second instance is the backward pass of a free monoid, and it is where the move does real work.

\subsubsection{The branching free monoid}
Read a container $C = (S \triangleleft P)$ bidirectionally: $S$ requests, $P s$ answers. Then the
composition product $C \triangleleft D$ is \emph{answer-dependent sequencing}---a first request $s$
together with, for each possible answer $p \in P s$, a second request---and the free monoid on $C$ under
$\triangleleft$ is the type of answer-driven interaction trees over $C$. This gadget is Hedges'
``Kleene star on containers'' \cite{Hedges2025}, which rebuilds a bidirectional typechecker on it but
proves neither its monoid laws nor a universal property.

Concretely \cite{GambinoKock2013}, the free monoid is $C^{*} = (S^{*} \triangleleft P^{*})$ where
$S^{*}$ is closed $P$-trees $t ::= \lf \mid \nd\, s\, \kappa$ with $s \in S$ and
$\kappa : P s \to S^{*}$, and $P^{*} t = \leaves(t)$ is the set of root-to-$\lf$ paths, recorded as the
tuple of positions taken. A leaf is thus a \emph{complete answer-history} of one run. The unit is
$\lf$; the multiplication grafts and splits,
\begin{align*}
  \graft(\lf,u) &= u(()), & \graft(\nd\, s\, \kappa, u) &= \nd\, s\,(\lambda p.\,\graft(\kappa_p, u_p)),\\
  \spl(\lf,u,z) &= \langle (), z\rangle, &
  \spl(\nd\, s\, \kappa, u, \langle p, z\rangle) &= \big\langle \langle p, \ell\rangle, w\big\rangle,
\end{align*}
where in the last line $(\ell,w) = \spl(\kappa_p, u_p, z)$ and $u_p \defeq \lambda r.\,u\langle p,r\rangle$.
The three monoid laws of $(C^{*},\graft,\lf)$ hold for every container; this is established elsewhere
(the universal property of the free monad on a container), and the underlying $\triangleleft$-monoid laws
are formalised in Lean. The backward pass of the multiplication is $\spl$.

\subsubsection{Path-confinement}
The worry Hedges' construction raises is the bidirectional one: since the continuation after a request
branches on the \emph{answer}, does that answer-dependence corrupt the backward pass $\spl$, so that it
fails to be natural in the grafted family and no universal property can hold? The following lemma says it
cannot, and is the content of the move in this setting.

\begin{lemma}[Path-confinement]\label{lem:D1}
For every container $C$, every tree $t \in S^{*}$, every family $u : \leaves(t) \to S^{*}$, and every
leaf $z \in \leaves(\graft(t,u))$,
\[
  \spl(t,u,z) = (\ell, w),
\]
where $\ell$ is the \emph{unique prefix of $z$ that is a leaf of $t$} and $w$ is the corresponding
suffix. In particular $\spl(t,u,z)$ is independent of $u$: the backward pass of the multiplication reads
only the outer tree.
\end{lemma}

\begin{proof}
Structural induction on $t$, using guise (4) of Lemma~\ref{lem:rootpath}.

\emph{Base $t = \lf$.} By definition $\spl(\lf,u,z) = \langle(),z\rangle$. The empty path $()$ is a prefix
of every $z$ and is the unique leaf of $\lf$, so $\ell = ()$, $w = z$; the output did not consult $u$.

\emph{Step $t = \nd\, s\, \kappa$.} A leaf of $\graft(\nd\, s\, \kappa, u) = \nd\, s\,(\lambda p.\,
\graft(\kappa_p, u_p))$ has the form $z = \langle p, z'\rangle$ with $p \in P s$ and
$z' \in \leaves(\graft(\kappa_p, u_p))$. By definition $\spl(\nd\, s\, \kappa, u, z) =
\langle\langle p, \ell'\rangle, w'\rangle$ with $(\ell', w') = \spl(\kappa_p, u_p, z')$. By the induction
hypothesis applied to the subtree $\kappa_p$ with family $u_p$, the pair $(\ell', w')$ depends only on
$\kappa_p$ and $z'$, and $\ell'$ is the unique $\kappa_p$-leaf prefix of $z'$. Now a prefix of
$z = \langle p, z'\rangle$ is a leaf of $\nd\, s\, \kappa$ iff it has the form $\langle p, r\rangle$ with
$r$ a leaf-prefix of $z'$ in $\kappa_p$; such $r$ exists and is unique, namely $r = \ell'$. Hence
$\langle p, \ell'\rangle$ is the unique $(\nd\, s\, \kappa)$-leaf prefix of $z$, with suffix $w'$, and the
computation used only $t$ and $z$.
\end{proof}

\begin{corollary}\label{cor:no-obstruction}
The answer-dependence of a container's continuations is confined to the shape data of $C^{*}$ (the trees
and their per-answer children) and to the leaves (the complete answer-histories); it never enters the
backward pass of the monoid multiplication. The one-directional proof of the universal property of
$C^{*}$ is therefore already the bidirectional proof, and the feared obstruction does not exist.
\end{corollary}

\begin{proof}
The backward component of associativity is, by Lemma~\ref{lem:D1}, a statement purely about paths: both
split-composites of the associator read only their respective outer trees, so the double split of a leaf
of $\graft(\graft(t,u),v)$ is the unique three-fold factorisation into outer-leaf, middle-leaf, and inner
suffix, whose well-definedness is associativity of path concatenation; no answer datum from $u$ or $v$ is
read. The unit laws are trivial. Hence no monoid law of $C^{*}$ consults the grafted data in its backward
component.
\end{proof}

\begin{remark}[The linear case is O'Neill]
When $S$ is a singleton the trees are lists and $\spl$ is list-splitting at a prefix, the inverse of
concatenation---exactly the backward law of O'Neill's tensor-algebra free $\triangleleft$-monoid. The
branching case runs the same list-splitting path by path: Hedges' Kleene star is the indexed
generalisation of O'Neill's tensor algebra, with branching supplied by the shapes and the backward law
unchanged.
\end{remark}

\emph{Status.} Lemma~\ref{lem:D1} is proved above and verified computationally on $144{,}882$ triples
$(t,u,z)$ over three containers, including a multi-shape bidirectional typechecker with a childless
(answer-free) request and a branch-deleting graft; the $\triangleleft$-monoid laws of $C^{*}$ are
formalised in Lean, while a Lean formalisation of Lemma~\ref{lem:D1} itself is a flagged target and not
yet done. The universal property it serves is prior work; the new content is the isolation of the
backward pass as path-confined.

\subsection{Instance (III): factoring through the least element}\label{sec:arboreal}
The third instance is an axiom in the theory of arboreal categories, and we cite it.

Arboreal categories were introduced by Abramsky and Shah \cite{AbramskyShah2021} as an axiomatic home for
the game comonads of finite model theory---the Ehrenfeucht--Fra\"iss\'e, pebbling, and modal comonads,
whose coalgebras encode plays of model-comparison games. A \emph{path} in such a category is an object
whose poset of subobjects (for a fixed proper factorisation system) is a finite linear order
\cite[Def.~4]{JaklReggio2026}; paths are the chains, the single branches. The original axioms required
paths to be \emph{connected}: every arrow from a path into a coproduct $\coprod_{i} Q_i$ of paths
factors through a single summand \cite[Def.~5]{JaklReggio2026}.

As Jakl and Reggio observe \cite[Ex.~8]{JaklReggio2026}, this fails for the modal comonads: at the level
of the underlying forest orders, coproducts glue trees at the root, and a path through such a gluing need
not land in one summand. Their repair replaces connectedness with \emph{tree-connectedness}.

\begin{definition}[tree-connectedness, {\cite[Def.~11]{JaklReggio2026}}]\label{def:tc}
A path $P$ is \emph{tree-connected} if for every tree-diagram of paths $D : T \to \A$ and every arrow
$f : P \to \operatorname{colim} D$ there is an element $t \in T$ through which $f$ factors via the colimit
coprojection, and $t$ is the \emph{least} such element of $T$.
\end{definition}

This is Schema~\ref{sch:main} in guise (3): the factoring is through the \emph{least} element of the tree
$T$, the root of the relevant down-set. Jakl and Reggio show that tree-connectedness recovers all the
properties that the connectedness axiom was used to prove, now including the modal comonads
\cite[Lem.~12]{JaklReggio2026}, and that the resulting path functor is a Street fibration whose cartesian
morphisms are exactly the \emph{pathwise embeddings} \cite[Thm.~23, Thm.~26]{JaklReggio2026}. We will use
these two theorems in Section~\ref{sec:open}; for now the point is only that the repair \emph{is} the
move: replace an arbitrary-coproduct condition by factoring through a least tree element.

\subsection{Instance (IV): the base is a forest}\label{sec:forest}
The fourth instance makes the move a property of a small category, and it is the one that becomes a
theorem.

\subsubsection{Coalgebras of a polynomial comonad}
\begin{lemma}\label{lem:coalg}
Let $c = C$ be a polynomial comonad, viewed as a small category via
Proposition~\ref{prop:comonoid-cat}. The category of Eilenberg--Moore coalgebras of the comonad
$\cext{c}$ on $\Set$ is equivalent to the copresheaf category $[C,\Set]$, equivalently to the discrete
opfibrations over $C$.
\end{lemma}

\begin{proof}
A coalgebra is a map $\gamma : X \to \cext{c}X = \sum_{s \in S} X^{P s}$. Write
$\gamma(x) = \big(|x|,\, (x \cdot f)_{f \in P|x|}\big)$ with $|x| \in S = \ob\, C$ the \emph{label} of $x$
and $x \cdot f \in X$. The counit law $\varepsilon \circ \gamma = \id$ says $x \cdot \id_{|x|} = x$. The
coassociativity law $(\delta \triangleleft \id)\circ\gamma = (\id \triangleleft \gamma)\circ\gamma$ says,
for $f : |x| \to a'$ in $C$ and $g : a' \to a''$, that $|x \cdot f| = \cod(f)$ and
$x \cdot (g \circ f) = (x \cdot f)\cdot g$. These are exactly the data and laws of a functor
$X : C \to \Set$ with $X(a) = \{x : |x| = a\}$ and $X(f)(x) = x \cdot f$: an object over each object, a chosen
lift of each arrow out of it, respecting identities and composition---i.e.\ a discrete opfibration over
$C$. Coalgebra morphisms are natural transformations.
\end{proof}

\subsubsection{Arboreal $=$ forest}
\begin{corollary}\label{cor:arboreal-forest}
If $C$ is a forest category then every object of $[C,\Set]$ is a forest, by Lemma~\ref{lem:rootpath}(1).
For $C = [k]$ the finite linear order this equivalence is exact: $[C,\Set]$ is the category of forests of
height $\le k$; for $C = [k]^{+}$ the prefix trie, it is the trie-indexed forests. Hence the arboreal
coalgebra categories of the Ehrenfeucht--Fra\"iss\'e and pebbling comonads are the copresheaf categories
of their (forest) comonoids: for these game comonads, ``\emph{arboreal $=$ the base is a forest}'' is a
theorem, not an analogy.
\end{corollary}

\begin{proof}
A copresheaf $X : [k] \to \Set$ assigns a set $X_n$ at each $n$ and, for each arrow $n \to m$ (which
exists iff $m \le n$), an ``ancestor'' map $X_n \to X_m$. Grade each element by its label and order
$x \le y$ iff $|x| \le |y|$ and $y \cdot |x| = x$; then every principal down-set
$\{y \cdot m : m \le |y|\}$ is a chain of length $|y| \le k$---a forest of height $\le k$ by
Lemma~\ref{lem:rootpath}(1)---and depth-grading a forest inverts the construction. The trie case is
identical with $[k]^{+}$ in place of $[k]$.
\end{proof}

By Proposition~\ref{prop:comonoid-cat} and a direct computation, the Ehrenfeucht--Fra\"iss\'e comonad
$\EF_k$ (nonempty sequences of length $\le k$, counit last, comultiplication prefixes) has comonoid the
linear order $[k]$, and the pebbling comonad $\Peb_k$ has comonoid the $k$-ary prefix trie $[k]^{+}$; the
modal comonad is not polynomial over $\Set$, consistently with its being the instance that forced the
tree-connectedness repair of Section~\ref{sec:arboreal}. Thus Corollary~\ref{cor:arboreal-forest}
identifies the arboreal coalgebra categories of $\EF_k$ and $\Peb_k$---the forests of bounded height and
the trie-indexed forests of Abramsky and Shah \cite{AbramskyShah2021}---with the copresheaf categories of
their polynomial comonoids. This is the bridge promised in the introduction.

\emph{Status and novelty.} Lemma~\ref{lem:coalg} is standard-adjacent (the familial / cofree-coalgebra
story in Poly \cite{SpivakGarnerFairbanks2025}) and Corollary~\ref{cor:arboreal-forest}'s forest
identification is proved directly for $\EF_k$ by Abramsky and Shah \cite{AbramskyShah2021}; the increment
claimed here is only the clean identification of the two, promoting ``arboreality $=$ forest'' to a
theorem and exhibiting it as the sharpest form of Schema~\ref{sch:main}. The coalgebra counts agree with
forest counts in every tested case ($|X|=2,3$ and $k=2,3$).

\section{Reading the four as one, and an open problem}\label{sec:open}

\subsection{One move, four sharpenings}
The four instances are the same move at four altitudes, each sharper than the last:
\[
  \underbrace{\text{root at each shape}}_{\text{(I) positions}}
  \;\longrightarrow\;
  \underbrace{\text{cut at the unique prefix}}_{\text{(II) free-monoid backward pass}}
  \;\longrightarrow\;
  \underbrace{\text{factor through the least element}}_{\text{(III) arboreal colimits}}
  \;\longrightarrow\;
  \underbrace{\text{the base category is a forest}}_{\text{(IV) polynomial comonoid}}.
\]
Instance (I) is a discipline on a position set; instance (II) a property of a backward map; instance (III)
a property of colimit factorisations; instance (IV) a property of a small category. Lemma~\ref{lem:rootpath}
is what they share: in each, the data below a node is a chain, so the operative map reads a unique
root-path and ignores all branching. That the move appears independently in a frontier paper on arboreal
categories \cite{JaklReggio2026} and in the backward pass of a free monoid, on structurally the same day
and with neither borrowing from the other, is the evidence that it is a genuine normal form and not a
coincidence of vocabulary.

\subsection{Where the schema would next bite: an open problem}
Instance (IV) identifies the arboreal coalgebra category with a copresheaf category \emph{as a category}.
The deeper question is whether the match persists \emph{as a fibration}. Jakl and Reggio's path functor
$\pathfun : \A \to \T$ is a Street fibration whose cartesian morphisms are the pathwise embeddings
\cite[Thm.~23, Thm.~26]{JaklReggio2026}; the Poly side has its own container fibration
$\Cont(\cod) = \Fam(\cod^{\op})$, a bifibration over the category of containers. Do these coincide?

\begin{conjecture}[refined crown; \emph{speculative}]\label{conj:crown}
For $c = C$ a forest polynomial comonad, the restriction of the arboreal path fibration
$\pathfun : \A \to \T$ to the subcategory of paths is equivalent, as a fibration, to a $\Fam$-type
fibration built from $C$ via $\Cont(\cod)$---up to the fibrewise opposite forced by the two sides
computing dual (co-)intuitionistic predicate structures.
\end{conjecture}

We state Conjecture~\ref{conj:crown} as \emph{speculative}, and we can say precisely how far the evidence
reaches and where it stops. On the \emph{unary, propositionally-truncated, single-path} fragment---one
unary relation, subobjects rather than proof-relevant predicates, a chain rather than a branching
tree---we have computed (by explicit enumeration at $C = [2],[3],[4]$ together with a $k$-independent
derivation) the following, graded \emph{computed}:

\begin{itemize}
\item the two fibres over a path $[n]$ are \emph{opposite} posets---$\pathfun$ gives colourings under
$\subseteq$, $\Cont(\cod)$ gives subobjects under $\supseteq$;
\item the arboreal cartesian reindexing along a path inclusion is \emph{preimage}
$R \mapsto R \cap [m]$---relation-reflection, substitution of predicates;
\item this preimage is the \emph{opcartesian} (pullback) leg $\rho^{*}$ of the bifibration
$\Cont(\cod)$, \emph{not} its cartesian leg $(\Sigma_\rho)^{\op}$, which is direct image and differs from
preimage already at the deepest singleton.
\end{itemize}

So on this fragment the arboreal path fibration is the fibrewise opposite of $\Cont(\cod)$, and the two
theories meet precisely on the pullback leg $\rho^{*}$: relation-reflection is $\rho^{*}$, not
$(\Sigma_\rho)^{\op}$. This is a usable dictionary---it tells any application that needs to ``pull a
predicate back along a sub-trajectory'' which leg of the container bifibration to use---but it is not the
conjecture. Three gaps keep Conjecture~\ref{conj:crown} open, and we name them honestly:

\begin{enumerate}
\item \emph{Path base only.} The single-shape container $\Phi([n]) = (1 \triangleleft [n])$ reaches only
the paths $\Paths \subset \T$, not the branching trees; the full base needs multi-shape containers whose
shapes are the branches, and the branch bookkeeping is not done.
\item \emph{Unary only.} A $k$-ary relation on $[n]$ is a subset of $[n]^{k}$, which is not an object of
$\Set/[n]$; higher-arity relations are invisible to $\Cont(\cod)$. The match is a statement about the
unary fragment, not about all of finite model theory's game comonads.
\item \emph{Truncation.} The comparison is at the level of subobjects. The proof-relevant
$\Cont(\cod)$-fibre $(\Set/[n])^{\op}$ is strictly richer than the arboreal $\Sub([n])$, and a
proof-relevant arboreal theory with witness-set colours is not standard.
\end{enumerate}

We flag, finally, that two of the founding game-comonad references \cite{AbramskyShah2021,
AbramskyDawarWang2017} are cited here through the survey \cite{JaklReggio2026} and should be read directly
before any formal submission; the load-bearing definitional content of this paper rests only on the
deep-read sources.

\section{Conclusion}\label{sec:conclusion}

One move---confine to a rooted tree, read the unique root-path---organises the root axiom of directed
containers, the backward pass of the branching free monoid, the tree-connectedness repair of arboreal
categories, and the forest condition on polynomial comonoids. Its single geometric content is that a
forest has no branching below a node (Lemma~\ref{lem:rootpath}), and in its sharpest form it is a theorem:
the coalgebras of a polynomial comonad $c = C$ are the copresheaves on $C$, so arboreality is exactly the
forest property of $C$ (Theorem~\ref{thm:main}). This identifies the arboreal coalgebra categories of
finite model theory with copresheaf categories of polynomial comonads---a bridge the two programmes had
not drawn.

Two directions stand out. The first is to settle whether the four instances are specialisations of one
\emph{lemma}---a single statement, over a small category $C$ with a factorisation system, equating
``connectivity'' with ``$C$ a forest'' and the operative map with the unique-prefix map---or whether they
remain a family united only in spirit; instance (IV) supplies the vocabulary to attempt it. The second is
Conjecture~\ref{conj:crown}: on the unary truncated path fragment the arboreal path fibration is the
fibrewise opposite of the container bifibration, meeting it on the pullback leg $\rho^{*}$; whether this
survives the passage to branching trees, higher arities, and proof-relevance is open, and the
obstruction in each case is named above. Either direction would turn a structural coincidence into a
structural theorem---which is, after all, the whole point of noticing that one keeps making the same move.

\section*{Acknowledgements}
The author thanks Neil Ghani and Robin Langer for the research setting, and Rick for the Bridge-3
material that prompted instance (IV).

\begin{thebibliography}{99}

\bibitem{AbramskyDawarWang2017}
S.~Abramsky, A.~Dawar, P.~Wang,
\emph{The pebbling comonad in finite model theory},
in: Proc.\ 32nd Annual ACM/IEEE Symposium on Logic in Computer Science (LICS), 2017.
arXiv:1704.05124.

\bibitem{AbramskyShah2021}
S.~Abramsky, N.~Shah,
\emph{Relating structure and power: comonadic semantics for computational resources},
Journal of Logic and Computation \textbf{31}(6) (2021), 1390--1428.
arXiv:1806.09031.

\bibitem{AhmanChapmanUustalu2014}
D.~Ahman, J.~Chapman, T.~Uustalu,
\emph{When is a container a comonad?},
Logical Methods in Computer Science \textbf{10}(3) (2014).

\bibitem{AhmanUustalu2016}
D.~Ahman, T.~Uustalu,
\emph{Directed containers as categories},
in: Proc.\ 6th Workshop on Mathematically Structured Functional Programming (MSFP), EPTCS 207, 2016,
89--98.

\bibitem{GambinoKock2013}
N.~Gambino, J.~Kock,
\emph{Polynomial functors and polynomial monads},
Mathematical Proceedings of the Cambridge Philosophical Society \textbf{154}(1) (2013), 153--192.

\bibitem{Hedges2025}
J.~Hedges,
\emph{Bidirectional typechecking with dependent lenses},
Cybernetics and Category Theory Institute (cybercat.institute), 24 February 2025.

\bibitem{JaklReggio2026}
T.~Jakl, L.~Reggio,
\emph{On the axioms of arboreal categories},
2026. arXiv:2603.21841.

\bibitem{SpivakGarnerFairbanks2025}
D.~I.~Spivak, R.~Garner, J.~Fairbanks,
\emph{Functorial aggregation},
Journal of Pure and Applied Algebra \textbf{229} (2025), 107883.

\end{thebibliography}

\end{document}
